% !TEX root =./ECE444_Practice_main.tex
%
% L7 practice set (Simple Resonant Antennas) -- the isotropic reference, EIRP,
% the half-wave dipole pattern and directivity, resonance and the VSWR that
% follows from it, and a design-at-frequency problem. Monopoles and loops are
% deliberately absent; they belong to L9.

\fullwidth{\textbf{Learning Objective 2.1: } I can describe the radiation behavior of simple
resonant antennas (isotropic radiator, half-wave dipole, monopole, loop) and calculate their
gain and impedance.
\\[4pt]\normalfont\small Show your work. A \textbf{Documentation} line at the end
is required (write \textbf{None} if you did not collaborate).}

\begin{questions}

%%%%%%%%%%%%%%%%%%%%%%%%%%%%%%%%%%%%%%%%%%%%%%%%%%%%%%%%%%%%%%%%%%%%%%%%%%%%
\question \textbf{The reference that cannot be built.} Every gain figure you will
ever quote is a ratio against an antenna nobody has ever made.

\begin{parts}

	\part Explain, in one or two sentences, why a truly isotropic radiator cannot exist.
	Your answer must mention what a real antenna is forced to have that an isotropic
	radiator does not.
		\begin{solution}[1.0in]
		In the far field the electric field is transverse, so it lies tangent to any sphere
		surrounding the antenna. A nonvanishing tangential vector field cannot be defined
		everywhere on a sphere (the ``hairy ball'' result), so the field must go to zero
		somewhere. That zero is a \textbf{null}. Every real antenna therefore has at least one
		null; an isotropic radiator by definition has none, so it is a unit of comparison
		rather than a buildable antenna.
		\end{solution}

	\begin{minipage}[t]{0.58\linewidth}
		\part A transmitter delivers $5\watts$ to a lossless half-wave dipole.
		Find the EIRP in watts and in \dbm.
		\begin{solution}[1.2in]
		The half-wave dipole has $D = 1.64$, and with negligible conductor loss $G = D$.
		\eq{
		\text{EIRP} &= P_t G_t = \lp 5\watts \rp \lp 1.64 \rp = 8.2\watts \\
		P_t \lb \dbm \rb &= 10\mylog{5000} = 37.0\dbm \\
		\text{EIRP} &= 37.0\dbm + 2.15\dbi = 39.1\dbm
		}
		\end{solution}
	\end{minipage}\hfill%
	\begin{minipage}[t]{0.37\linewidth}
		\raggedleft \vspace{12pt}
		EIRP = \ansbox[1.5in]{$8.2\watts = 39.1\dbm$}
	\end{minipage}

	\begin{minipage}[t]{0.58\linewidth}
		\part A vendor lists an antenna at $8\db$ over a dipole. What transmit power is
		needed to reach an EIRP of $40\dbm$?
		\begin{solution}[1.2in]
		Convert the gain to the isotropic reference first:
		\eq{
		G &= 8 \text{ dBd} + 2.15 = 10.15\dbi \\
		P_t &= \text{EIRP} - G = 40\dbm - 10.15\dbi = 29.9\dbm \\
		P_t &= 10^{29.85/10}\mwatts = 966\mwatts \approx 0.97\watts
		}
		\end{solution}
	\end{minipage}\hfill%
	\begin{minipage}[t]{0.37\linewidth}
		\raggedleft \vspace{12pt}
		$P_t$ = \ansbox[1.5in]{$29.9\dbm = 0.97\watts$}
	\end{minipage}

	\part Why do spectrum regulators write emission limits in EIRP rather than in
	transmitter power? Answer in one or two sentences.
		\begin{solution}[0.9in]
		Interference depends on the power density arriving at a victim receiver, which is set
		by the transmitter power \emph{and} the antenna gain together. EIRP combines the two
		into the single quantity that actually determines the field in a given direction, so a
		limit on EIRP cannot be evaded by keeping the transmitter small and bolting on a very
		high-gain antenna.
		\end{solution}

\end{parts}

\newpage
%%%%%%%%%%%%%%%%%%%%%%%%%%%%%%%%%%%%%%%%%%%%%%%%%%%%%%%%%%%%%%%%%%%%%%%%%%%%
\question \textbf{Cut a dipole for the 915 MHz ISM band.} You are building a resonant
half-wave dipole for a $915\mhz$ telemetry link out of thin copper wire.

\begin{parts}

	\begin{minipage}[t]{0.58\linewidth}
		\part Find the free-space wavelength, the resonant total length using the 5\% rule,
		and the length of each arm.
		\begin{solution}[1.4in]
		\eq{
		\lambda &= \frac{c}{f} = \frac{3\ee{8}}{915\ee{6}} = 0.328\m = 32.8\text{cm} \\
		L &= 0.475\lambda = 0.475 \lp 32.8\text{cm} \rp = 15.6\text{cm} \\
		\frac{L}{2} &= 7.8\text{cm} \text{ per arm}
		}
		Cross-check with the field form: $143/915 = 0.156\m$, the same answer.
		\end{solution}
	\end{minipage}\hfill%
	\begin{minipage}[t]{0.37\linewidth}
		\raggedleft \vspace{12pt}
		$L$ = \ansbox[1.3in]{$15.6\text{cm}$}
		\\[10pt]
		arm = \ansbox[1.3in]{$7.8\text{cm}$}
	\end{minipage}

	\begin{minipage}[t]{0.58\linewidth}
		\part Compute $2D^2/\lambda$ for this antenna. Then state how far away you would
		actually stand to measure its pattern, and why that is not the same number.
		\begin{solution}[1.4in]
		\eq{
		\frac{2D^2}{\lambda} &= \frac{2\lp 0.156\m \rp^2}{0.328\m} = 0.148\m = 14.8\text{cm}
		}
		That is less than half a wavelength, so the far-field criterion from L5 is not the
		binding one here. For an electrically small antenna the reactive near field is the
		limitation, and the practical rule is $r \gg \lambda$ --- stand at several
		wavelengths, say $5\lambda \approx 1.6\m$ or more. The $2D^2/\lambda$ criterion only
		takes over once the antenna is large compared with a wavelength.
		\end{solution}
	\end{minipage}\hfill%
	\begin{minipage}[t]{0.37\linewidth}
		\raggedleft \vspace{12pt}
		$2D^2/\lambda$ = \ansbox[1.3in]{$14.8\text{cm}$}
		\\[10pt]
		use $r >$ \ansbox[1.3in]{$\approx 1.6\m$}
	\end{minipage}

	\begin{minipage}[t]{0.58\linewidth}
		\part The radio delivers $100\mwatts$ to this dipole. Find the EIRP in \dbm.
		\begin{solution}[1.0in]
		\eq{
		P_t &= 10\mylog{100} = 20.0\dbm \\
		\text{EIRP} &= 20.0\dbm + 2.15\dbi = 22.2\dbm = 164\mwatts
		}
		\end{solution}
	\end{minipage}\hfill%
	\begin{minipage}[t]{0.37\linewidth}
		\raggedleft \vspace{12pt}
		EIRP = \ansbox[1.4in]{$22.2\dbm$}
	\end{minipage}

	\begin{minipage}[t]{0.58\linewidth}
		\part You build it and the analyzer shows resonance at $880\mhz$, not $915\mhz$.
		Which way must you trim, and by how much per arm? Justify the direction in words.
		\begin{solution}[1.4in]
		Resonance low means the antenna is electrically \textbf{too long} --- a longer wire
		resonates at a lower frequency --- so trim it shorter. Length scales inversely with
		the resonant frequency:
		\eq{
		L_\text{new} &= L_\text{old}\frac{f_\text{old}}{f_\text{new}} = 15.6\text{cm} \lp \frac{880}{915} \rp = 15.0\text{cm} \\
		\Delta L &= 0.6\text{cm} \text{ total}, \quad 0.3\text{cm} \text{ per arm}
		}
		Trim a little less than that and re-measure; you cannot put wire back.
		\end{solution}
	\end{minipage}\hfill%
	\begin{minipage}[t]{0.37\linewidth}
		\raggedleft \vspace{12pt}
		trim \ansbox[1.4in]{$0.3\text{cm}$ per arm}
	\end{minipage}

\end{parts}

\newpage
%%%%%%%%%%%%%%%%%%%%%%%%%%%%%%%%%%%%%%%%%%%%%%%%%%%%%%%%%%%%%%%%%%%%%%%%%%%%
\question \textbf{Reading the half-wave pattern.} The normalized field pattern of a
half-wave dipole lying along $z$ is
\eq{
\left| F(\theta) \right| = \frac{\cosp{\frac{\pi}{2}\cos\theta}}{\sin\theta}
}

\begin{parts}

	\begin{minipage}[t]{0.58\linewidth}
		\part Evaluate the pattern at $\theta = 45\mydeg$ and express it in \db\ relative
		to the broadside peak.
		\begin{solution}[1.3in]
		\eq{
		\cos 45\mydeg &= 0.7071, \quad \frac{\pi}{2}\lp 0.7071 \rp = 1.1107 \text{ rad} \\
		\left| F \right| &= \frac{\cosp{1.1107}}{\sin 45\mydeg} = \frac{0.4441}{0.7071} = 0.628 \\
		\left| F \right| \lb \db \rb &= 20\mylog{0.628} = -4.0\db
		}
		\end{solution}
	\end{minipage}\hfill%
	\begin{minipage}[t]{0.37\linewidth}
		\raggedleft \vspace{12pt}
		$\left| F \right|$ = \ansbox[1.4in]{$0.628 = -4.0\db$}
	\end{minipage}

	\begin{minipage}[t]{0.58\linewidth}
		\part Show that $\theta = 51\mydeg$ is a half-power point, and use it to state
		the HPBW.
		\begin{solution}[1.3in]
		\eq{
		\cos 51\mydeg &= 0.6293, \quad \frac{\pi}{2}\lp 0.6293 \rp = 0.9886 \text{ rad} \\
		\left| F \right| &= \frac{0.5502}{0.7771} = 0.708 \approx \frac{1}{\sqrt{2}} \\
		\left| F \right| \lb \db \rb &= 20\mylog{0.708} = -3.0\db
		}
		By symmetry about broadside the other half-power point is at $129\mydeg$, so
		$\theta_\text{HP} = 129\mydeg - 51\mydeg = 78\mydeg$.
		\end{solution}
	\end{minipage}\hfill%
	\begin{minipage}[t]{0.37\linewidth}
		\raggedleft \vspace{12pt}
		$\theta_\text{HP}$ = \ansbox[1.3in]{$78\mydeg$}
	\end{minipage}

	\begin{minipage}[t]{0.58\linewidth}
		\part The half-wave dipole has $D = 1.64$. Find its directivity in the direction
		$\theta = 45\mydeg$, in \dbi, and comment on the sign.
		\begin{solution}[1.3in]
		Directivity in a given direction is the peak value scaled by the normalized power
		pattern:
		\eq{
		D(45\mydeg) &= D \left| F(45\mydeg) \right|^2 = 1.64 \lp 0.628 \rp^2 = 0.647 \\
		D(45\mydeg) \lb \dbi \rb &= 10\mylog{0.647} = -1.9\dbi
		}
		It is \emph{negative}, meaning the dipole radiates less in that direction than an
		isotropic radiator would. A directivity above 1 in one direction is always paid for
		with a directivity below 1 somewhere else.
		\end{solution}
	\end{minipage}\hfill%
	\begin{minipage}[t]{0.37\linewidth}
		\raggedleft \vspace{12pt}
		$D(45\mydeg)$ = \ansbox[1.3in]{$-1.9\dbi$}
	\end{minipage}

	\part The polar plot of a half-wave dipole looks strongly directional, yet its
	directivity is only $1.64$ ($2.15\dbi$). Explain the apparent contradiction in one or
	two sentences.
		\begin{solution}[1.0in]
		The plot is a single cut in a plane containing the wire. The pattern has \emph{no}
		$\phi$ dependence at all, so the antenna radiates equally in every azimuth direction
		--- it is a doughnut, not a beam. Directivity integrates over the whole sphere, and
		concentrating power in a $78\mydeg$-wide band that wraps the full $360\mydeg$ of
		azimuth buys very little concentration overall.
		\end{solution}

\end{parts}

\newpage
%%%%%%%%%%%%%%%%%%%%%%%%%%%%%%%%%%%%%%%%%%%%%%%%%%%%%%%%%%%%%%%%%%%%%%%%%%%%
\question \textbf{Impedance, resonance, and the match.} A thin dipole cut to exactly
$\lambda/2$ presents an input impedance of $73 + j42.5$\ohms; trimmed to resonance it
presents about $70 + j0$\ohms at its terminals.

\begin{parts}

	\begin{minipage}[t]{0.58\linewidth}
		\part Find $\left| \Gamma \right|$ and the VSWR for the untrimmed dipole on a
		$50$\ohms line.
		\begin{solution}[1.4in]
		\eq{
		\Gamma &= \frac{Z_{in} - Z_0}{Z_{in} + Z_0} = \frac{23 + j42.5}{123 + j42.5} \\
		\left| \Gamma \right| &= \frac{\sqrt{23^2 + 42.5^2}}{\sqrt{123^2 + 42.5^2}} = \frac{48.3}{130.1} = 0.371 \\
		\text{VSWR} &= \frac{1 + 0.371}{1 - 0.371} = 2.18
		}
		\end{solution}
	\end{minipage}\hfill%
	\begin{minipage}[t]{0.37\linewidth}
		\raggedleft \vspace{12pt}
		VSWR = \ansbox[1.3in]{$2.18$}
	\end{minipage}

	\begin{minipage}[t]{0.58\linewidth}
		\part Repeat for the trimmed, resonant dipole on a $50$\ohms line, and again on a
		$75$\ohms line.
		\begin{solution}[1.4in]
		\eq{
		\left| \Gamma_{50} \right| &= \left| \frac{70 - 50}{70 + 50} \right| = \frac{20}{120} = 0.167 \\
		\text{VSWR}_{50} &= \frac{1 + 0.167}{1 - 0.167} = 1.40 \\
		\left| \Gamma_{75} \right| &= \left| \frac{70 - 75}{70 + 75} \right| = \frac{5}{145} = 0.034 \\
		\text{VSWR}_{75} &= \frac{1 + 0.034}{1 - 0.034} = 1.07
		}
		\end{solution}
	\end{minipage}\hfill%
	\begin{minipage}[t]{0.37\linewidth}
		\raggedleft \vspace{12pt}
		VSWR$_{50}$ = \ansbox[1.2in]{$1.40$}
		\\[10pt]
		VSWR$_{75}$ = \ansbox[1.2in]{$1.07$}
	\end{minipage}

	\begin{minipage}[t]{0.58\linewidth}
		\part What percentage of the incident power is reflected in the untrimmed
		$50$\ohms case, and in the trimmed $50$\ohms case?
		\begin{solution}[1.1in]
		\eq{
		\left| \Gamma \right|^2_\text{untrimmed} &= \lp 0.371 \rp^2 = 0.138 = 13.8\% \\
		\left| \Gamma \right|^2_\text{trimmed} &= \lp 0.167 \rp^2 = 0.028 = 2.8\%
		}
		\end{solution}
	\end{minipage}\hfill%
	\begin{minipage}[t]{0.37\linewidth}
		\raggedleft \vspace{12pt}
		$\left| \Gamma \right|^2$ = \ansbox[1.5in]{$13.8\%$ vs $2.8\%$}
	\end{minipage}

	\part Using your answers above, argue in one or two sentences whether an engineer with
	a $50$\ohms radio should spend effort trimming the dipole to resonance or sourcing
	$75$\ohms cable. Support the recommendation with the numbers.
		\begin{solution}[1.1in]
		Trimming. Removing the $+j42.5$\ohms of reactance takes the $50$\ohms VSWR from
		$2.18$ to $1.40$ --- from 13.8\% reflected power down to 2.8\% --- while the
		remaining $70$ versus $50$\ohms resistance mismatch is worth only about $0.4$ of a
		VSWR point. The reactance costs far more than the resistance offset does, and it is
		removed for free with a pair of side cutters.
		\end{solution}

\end{parts}

\newpage
%%%%%%%%%%%%%%%%%%%%%%%%%%%%%%%%%%%%%%%%%%%%%%%%%%%%%%%%%%%%%%%%%%%%%%%%%%%%
\question \textbf{Why the wire is never exactly half a wavelength.} These parts are
answered in words and with rules of thumb, not with long derivations.

\begin{parts}

	\part Explain physically why a resonant dipole is cut shorter than $\lambda/2$.
	Name the two mechanisms.
		\begin{solution}[1.2in]
		\textbf{End effect}: capacitance from the tips of the dipole to the surroundings lets
		charge accumulate past where the metal stops, so the standing wave behaves as though
		the wire extended further than it physically does. \textbf{Wire thickness}: a fatter
		element has more end capacitance and a lower characteristic impedance, shortening it
		further. Both mechanisms slow the wave traveling on the wire relative to free space,
		and a slower wave needs less physical length to make up an electrical half wavelength.
		Practical resonance therefore lands near $0.47\lambda$ to $0.48\lambda$.
		\end{solution}

	\part A dipole built from thick aluminum tubing is measured to resonate at
	$0.46\lambda$, shorter than the $0.475\lambda$ rule of thumb predicts. Is the
	measurement suspicious? Explain.
		\begin{solution}[1.0in]
		No --- it is exactly what the thickness mechanism predicts. The 5\% rule assumes
		ordinary thin wire; a fat tubular element carries more end capacitance and shortens
		further, so $0.46\lambda$ is a reasonable resonant length for it. The rule of thumb is
		a starting point for cutting, not a specification.
		\end{solution}

	\part A dipole one full wavelength long has a directivity of $3.82\dbi$, well above the
	half-wave dipole's $2.15\dbi$. Give the reason engineers do not simply use the longer
	one everywhere.
		\begin{solution}[1.1in]
		Its feed point is the problem, not its pattern. At $L = \lambda$ the standing-wave
		current has a \emph{null} at the center, so the terminals see a very small current for
		a given voltage and the input impedance climbs to hundreds or thousands of ohms ---
		wildly mismatched to any ordinary line and extremely sensitive to construction
		details. The half-wave dipole puts the current maximum at the feed instead, which is
		worth far more in practice than $1.7\db$ of directivity.
		\end{solution}

	\part At roughly what length does the broadside directivity of a straight dipole peak,
	what value does it reach, and what happens to the pattern beyond that length? State the
	design consequence in one sentence.
		\begin{solution}[1.2in]
		Broadside directivity peaks near $L \approx 1.25\lambda$ at about $5.2\dbi$. Past that
		length the accumulated phase reversals on the wire split the main lobe, and the peaks
		walk off broadside entirely --- by $1.5\lambda$ the strongest lobes sit near
		$43\mydeg$ and $137\mydeg$ and broadside gain has collapsed.
		\textbf{Consequence:} a straight wire cannot be made into a high-gain broadside
		antenna simply by making it longer; beyond about $1.25\lambda$ you must switch to an
		array, which is where Module 3 begins.
		\end{solution}

\end{parts}

\newpage
%%%%%%%%%%%%%%%%%%%%%%%%%%%%%%%%%%%%%%%%%%%%%%%%%%%%%%%%%%%%%%%%%%%%%%%%%%%%
\question \textbf{Using the machinery.} The radiation integral gave you a pattern
that works for \emph{any} dipole length,
\eq{
\left| F(\theta) \right| \propto \frac{\cosp{\frac{kL}{2}\cos\theta} - \cosp{\frac{kL}{2}}}{\sin\theta}
}
and the impedance came from two tabulated special functions, $C_{in}(2\pi) = 2.4376$ and
$Si(2\pi) = 1.4182$. Use them.

\begin{parts}

	\begin{minipage}[t]{0.58\linewidth}
		\part For a \textbf{full-wave} dipole ($L = \lambda$), evaluate the pattern at
		$\theta = 90\mydeg$ and $\theta = 60\mydeg$, then express the $60\mydeg$ value in
		\db\ relative to the peak.
		\begin{solution}[1.5in]
		With $L = \lambda$, $kL/2 = \pi$, so $\cosp{kL/2} = -1$.
		\eq{
		\theta = 90\mydeg: \quad \left| F \right| &= \frac{\cosp{0} - (-1)}{1} = 2 \\
		\theta = 60\mydeg: \quad \left| F \right| &= \frac{\cosp{\pi/2} - (-1)}{\sin 60\mydeg} = \frac{1}{0.8660} = 1.155 \\
		\text{normalized} &= \frac{1.155}{2} = 0.577 \\
		&= 20\mylog{0.577} = -4.8\db
		}
		The peak is still at $90\mydeg$, so a full-wave dipole is still broadside.
		\end{solution}
	\end{minipage}\hfill%
	\begin{minipage}[t]{0.37\linewidth}
		\raggedleft \vspace{12pt}
		peak at \ansbox[1.2in]{$90\mydeg$}
		\\[10pt]
		$60\mydeg$ value = \ansbox[1.2in]{$-4.8\db$}
	\end{minipage}

	\begin{minipage}[t]{0.58\linewidth}
		\part Compute the reactance of a half-wave dipole from $X = \frac{\eta}{4\pi}Si(2\pi)$,
		and state in one sentence why the wire radius does not appear in the answer.
		\begin{solution}[1.4in]
		\eq{
		\frac{\eta}{4\pi} &= \frac{377}{4\pi} = 30.0\ohms \\
		X &= 30.0 \lp 1.4182 \rp = 42.5\ohms
		}
		Every radius-dependent term in the induced-EMF expression is multiplied by
		$\sinp{kL}$, and at exactly $L = \lambda/2$ we have $kL = \pi$, so
		$\sinp{\pi} = 0$ and those terms vanish. The radius independence holds
		\emph{only} at exactly a half wavelength.
		\end{solution}
	\end{minipage}\hfill%
	\begin{minipage}[t]{0.37\linewidth}
		\raggedleft \vspace{12pt}
		$X$ = \ansbox[1.3in]{$42.5$\ohms}
	\end{minipage}

	\part The resistance came out of an integral of the far-field pattern over a sphere,
	but the reactance did not and could not. Explain why in one or two sentences.
		\begin{solution}[1.0in]
		A far-field power integral accounts only for power that \emph{leaves} the antenna and
		never returns, which is precisely the definition of the real part. Reactance
		represents energy stored in the reactive near field and handed back to the source
		every cycle; it never crosses the far-field sphere, so no amount of pattern
		integration can reveal it. Getting it requires a near-field calculation --- the
		induced-EMF method.
		\end{solution}

\end{parts}

\vspace{1.5em}
\noindent\textbf{Documentation:} \rule[-2pt]{0.6\linewidth}{0.4pt}

\end{questions}
